\scanpage{103}
% archive sheet number: 102
% source line 01
La partie homogène de $q$ est
% source line 02
\[
q_0(x,y,z)=x^2+y^2+z^2-\alpha yz+zx+xy\tag*{(19)}
\]
% source line 03
dont la matrice est
% source lines 04-06
\[
\deleted{\unsure{\widetilde q_0|\overrightarrow{q_0}}}
\begin{pmatrix}2&1&1\\1&2&-\alpha\\1&-\alpha&2\end{pmatrix}\tag*{(20)}
\]
% source lines 07-08
de discriminant
\[
\begin{aligned}
\delta&=\underline{8}-2\alpha-\underline{2}-2\alpha^2-\underline{2}\\
&=4-2(\underbrace{\alpha^2+\alpha}_{1})=2
\end{aligned}
\]
% source line 09
\[
\delta=2\qquad\text{donc}\qquad\underline{\delta'=1}\tag*{(21)}
\]
% source line 10
Forme quadratique donc \uuline{partout}, \uline{inversible}!

% source lines 11-12
Quant au discriminant de $q$ homogénéisé par une variable $t$
% source line 13
\[
q(x,y,z,t)=x^2+y^2+z^2-\alpha yz+zx+xy-t(x+y+z)
\]
% source line 14
c'est le déterminant de la matrice
% source lines 15-18
\[
\unsure{\widetilde q|\overrightarrow q}
\begin{pmatrix}
2&1&1&-1\\
1&2&-\alpha&-1\\
-\alpha&-1&2&-1\\
-1&-1&-1&0
\end{pmatrix}\tag*{(22)}
\]
% source lines 19-22; darker correction and additions at the foot
dont le discriminant \deleted{(sauf erreur de calcul)} est
$\deleted{\boxed{\alpha'}}\inserted{-\alpha-3}$,
\deleted{donc aussi \unsure{inv.} partout. Donc l'icosaèdre est inscrit}
\inserted{$\unsure{N|\mathcal N}(-\alpha-3)=5$ i.e. sauf en car. $5$ (\uline{NB} en car. $5$, l'icosaèdre est \unsure{inscrit} dans une sphère de rayon nul, donc quadrique \uline{sing.}!)}
